\begin{center}
\textsc{Gauss} an \textsc{Olbers}.\quad Göttingen, 18.~April 1822.
\end{center}

\par\noindent\hspace*{\parindent}.\ .\ .\ .\ .\ .\quad Bei allen meinen Rechnungen liegen folgende Dimensionen der
Erde zugrunde
\[
\begin{array}{@{}r@{\;}c@{\;}l@{\;}c@{}}
a & = & 3\,271\,821 & \text{[Toisen]} \\
b & = & 3\,261\,011 & \text{»} \\
\text{Abpl.} & = & \tfrac{1}{302{,}68}.\;\text{[\textsuperscript{*})]} &
\end{array}
\]
Eigentlich hatte ich ganz \textsc{Walbecks} Resultat annehmen wollen: \(\frac{1}{302{,}78}\).
Durch einen Schreibfehler hatte ich aber jene schon vor längerer Zeit der Berechnung
von mancherlei Hilfstafeln untergelegt, und hielt es um so weniger der Mühe
wert, diese deshalb umzuarbeiten, da der Unterschied weit unter der durch
alle Gradmessungen zu erreichenden Genauigkeit liegt.
